\section{Reconstructing targets from overlapping records}
\label{atl:section}

The atlas idea becomes precise when compatibility and information preservation
are treated as separate requirements. Local records should agree about the
features they share. Their combined information should also determine every
distinction required by the declared target. Agreement alone does not establish
the second property. This section supplies an exact linear model of both
requirements, a construction of minimal selective residuals, and a stability
bound for reconstruction. These are mathematical components proposed for this
architecture; they are not claims that an existing Hodge implementation already
provides the desired neural guarantee.

\subsection{Coordinate charts and partial views}
\label{atl:charts}

A genuine coordinate chart $\phi_i:U_i\to V_i$ is invertible on its domain.
Two charts describing the same underlying object have the transition map
\begin{equation}
  \tau_{ij}=\phi_j\circ\phi_i^{-1}
  \quad\text{on }\phi_i(U_i\cap U_j).
  \label{atl:transition}
\end{equation}
The identities
\begin{equation}
  \tau_{ji}=\tau_{ij}^{-1},
  \qquad
  \tau_{jk}\circ\tau_{ij}=\tau_{ik}
  \label{atl:cocycle}
\end{equation}
hold wherever the compositions are defined; the second identity is evaluated
on the common triple overlap. For affine transitions
$\tau_{ij}(u)=A_{ij}u+b_{ij}$, these conditions imply
\begin{align}
  A_{ji}A_{ij}&=I,
  & b_{ji}&=-A_{ji}b_{ij}, \nonumber\\
  A_{ik}&=A_{jk}A_{ij},
  & b_{ik}&=A_{jk}b_{ij}+b_{jk}.
  \label{atl:affine-cocycle}
\end{align}
The displayed matrix equalities apply when the affine maps agree on overlaps
with nonempty interior in their coordinate spaces. On a smaller constrained
overlap, equality is required only on that overlap.

A many-to-one layer is not automatically such a chart. For partial layer views,
use observation maps $r_i:\mathcal Z\to E_i$ and comparison maps
$\rho_{i,ij}:E_i\to O_{ij}$. The common overlap condition is
\begin{equation}
  \rho_{i,ij}(r_i(z))=\rho_{j,ij}(r_j(z)).
  \label{atl:restriction}
\end{equation}
Neither observation map needs an inverse. Their joint records may still
determine the task. Cellular sheaves provide a standard language for compatible
local data and restriction maps; their degree-zero global sections are
characterized by an agreement operator. The associated spectral theory also
studies reducing edge-message dimensions while preserving global sections
\citep{sheaves}. The proposal uses this organizational idea without assuming
that an arbitrary stack of neural layers is already an atlas or a sheaf with
the required observability properties.

Layer identities, overlap domains, branch choices, and directed transition
maps belong to the retained representation. Ordered composition in
\eqref{atl:cocycle} cannot generally be replaced by an unlabeled weighted sum.

\subsection{The source map and the agreement operator}
\label{atl:two-operators}

Let the source state be $z\in\R^n$, with local records $e_i=M_i z$.
Stack them into $e=Mz\in E=\bigoplus_i E_i$. Let
$\delta:E\to O$ collect overlap discrepancies, for example
$(\delta e)_{ij}=P_{i,ij}e_i-P_{j,ij}e_j$. Correct restrictions satisfy
\begin{equation}
  \delta M=0,
  \qquad\text{hence}\qquad
  \im M\subseteq\ker\delta.
  \label{atl:complex}
\end{equation}
This says that genuine source records agree. The stronger condition
\begin{equation}
  \im M=\ker\delta
  \label{atl:exactness}
\end{equation}
says that every compatible collection has an explanation within the chosen
source model. It is needed for that universal gluing claim, although decoding
records already known to equal $Mz$ does not require it.

The ambiguity left about the source is $\ker M$, not $\ker\delta$.
The latter contains compatible records themselves. In particular,
$\ker(\delta^\top\delta)=\ker\delta$ characterizes agreement; it does not
establish that those records determine the original input or target. If
\eqref{atl:exactness} holds, compatible records identify a source equivalence
class in $\R^n/\ker M$. The remaining question is whether the target is
constant on that class.

\subsection{An exact criterion for target observability}
\label{atl:observability}

Write the required linear task features as $t=Cz\in\R^q$. The rows of $C$
may describe several required readouts. Taking $C=I$ recovers the stronger
requirement of preserving the entire source.

\begin{theorem}[Linear target observability]
\label{atl:observability-theorem}
For unrestricted $z\in\R^n$, the following statements are equivalent:
\begin{enumerate}
  \item Some decoder recovers $Cz$ from $Mz$ for every $z$.
  \item $\ker M\subseteq\ker C$.
  \item $\row(C)\subseteq\row(M)$.
  \item There is a linear matrix $D$ such that $C=DM$.
\end{enumerate}
\end{theorem}

\begin{proof}
If a decoder exists and $v\in\ker M$, then $Mv=M0$, so it must give
$Cv=C0=0$. This proves the second statement. Conversely, that statement says
that two states with the same records have the same target. In finite
dimensional Euclidean spaces, $\row(M)=(\ker M)^\perp$ and likewise for $C$.
Taking orthogonal complements therefore makes the second statement equivalent
to the third. Under the third statement, express each row of $C$ as a linear
combination of rows of $M$ and collect the coefficients into $D$. This gives
$C=DM$. Finally, $D(Mz)=Cz$ provides the required decoder, completing all
implications.
\end{proof}

\begin{corollary}[A well-defined target for every compatible record]
\label{atl:gluing-corollary}
Assume $\im M=\ker\delta$. A linear gluing readout
$\Gamma:\ker\delta\to\R^q$ satisfying $\Gamma(Mz)=Cz$ exists if and only
if $\ker M\subseteq\ker C$. When it exists, one such formula is
$\Gamma(e)=CM^\dagger e$.
\end{corollary}

\begin{proof}
Exactness ensures that every compatible $e$ has a preimage $z$. Any two
preimages differ by a member of $\ker M$. Under the stated kernel inclusion,
they have the same target, so assigning $\Gamma(e)=Cz$ is well-defined.
Preimages can be added and scaled, which makes this assignment linear.
Moreover, $CM^\dagger M=C$ by target observability, proving the formula.
Conversely, the zero record has both $0$ and every $v\in\ker M$ as
preimages. A well-defined readout therefore requires $Cv=0$ for every such
$v$, establishing necessity.
\end{proof}

Thus individual views can be lossy while their combined records remain exactly
task-sufficient. For a restricted source domain $\mathcal Z$, the corresponding
necessary and sufficient condition is
\begin{equation}
  v\in\ker M\cap(\mathcal Z-\mathcal Z)
  \quad\Longrightarrow\quad Cv=0.
  \label{atl:restricted-domain}
\end{equation}
The unrestricted theorem is a convenient stronger certificate. It is also
necessary when the allowed source domain contains an open region: any missing
direction can then be scaled into a difference between allowed states.

\Needspace{20\baselineskip}
\subsection{The smallest selective linear residual}
\label{atl:minimal-residual}

Suppose $Mz$ is already retained. Additional records $r=Nz$ should supply the
target distinctions absent from it. The following result gives their minimum
dimension and an explicit decoder. Here $M^\dagger$ denotes the Moore--Penrose
pseudoinverse.

\begin{theorem}[Minimum residual dimension]
\label{atl:residual-theorem}
The minimum number of additional independent real-valued linear coordinates
needed to recover $Cz$ from $(Mz,Nz)$ for every $z$ is
\begin{equation}
  r_*=
  \rank\!\begin{bmatrix}M\\C\end{bmatrix}-\rank M
  =\rank(CP_K),
  \qquad P_K=I-M^\dagger M.
  \label{atl:minimum-rank}
\end{equation}
Choose $N$ with orthonormal rows spanning $\row(CP_K)$. Then $N$ has $r_*$
rows and the exact decoder is
\begin{equation}
  Cz=C M^\dagger(Mz)+C N^\top(Nz).
  \label{atl:residual-decoder}
\end{equation}
For $r_*=0$, the second term is absent.
\end{theorem}

\begin{proof}
The matrix $P_K$ is the orthogonal projector onto $K=\ker M$.
For any successful residual $N$, its restriction to $K$ must distinguish
everything distinguished there by $C$. The observability theorem, applied on
$K$, gives a factorization of $C|_K$ through $N|_K$. Consequently
\begin{equation*}
  \rank(CP_K)=\rank(C|_K)
  \le\rank(N|_K)\le\text{number of rows of }N.
\end{equation*}
This is the lower bound.

For the proposed construction, the rows of $N$ lie in $K$, so $NP_K=N$.
The matrix $N^\top N$ projects onto $\row(CP_K)$, giving
$CP_K=CP_KN^\top N=CN^\top N$. Therefore
\begin{equation*}
  C=C(I-P_K)+CP_K=CM^\dagger M+CN^\top N,
\end{equation*}
which proves the decoder and attains the bound. Finally,
$\row\bigl([M^\top,C^\top]^\top\bigr)$ is the orthogonal direct sum of
$\row(M)$ and $\row(CP_K)$: subtracting $CM^\dagger M$ from $C$ produces
the second summand, and adding it back recovers the original rows. Their
dimensions establish the rank-difference identity.
\end{proof}

This is a precise sense in which residual retention can be selective.
For $C=I$, $r_*=n-\rank M$; a smaller task family may need fewer coordinates.
The result counts real coordinates, not bits or computational cost. Dense
projections may be expensive, and near-zero singular values need a numerical
stability analysis even when symbolic rank is full.

The construction also requires a declared target map. If $C$ is known only
through examples, a learned estimate of its missing directions does not inherit
an all-input guarantee merely by using this formula. One controlled experiment
can supply $C$ analytically and verify the construction; a second can learn it
and measure the resulting errors. These are different claims. To preserve
several linear targets, stack their matrices into $C$ before constructing $N$.
Adding a future target can increase $\rank(CP_K)$ and require another retained
coordinate. An apparently unnecessary direction may therefore be dispensable
for one fixed readout yet necessary for reusable computation. The selected task
family must remain part of the specification rather than an implicit property
of the current training batch.

\subsection{A three-variable example}
\label{atl:example}

Let $z=(a,b,c)^\top$ and retain two adjacent measurements,
\begin{equation}
  s_{12}=a+b,\qquad s_{23}=b+c,
  \qquad
  M=\begin{bmatrix}1&1&0\\0&1&1\end{bmatrix}.
  \label{atl:adjacent-example}
\end{equation}
Its kernel is $\operatorname{span}\{(1,-1,1)^\top\}$. Individual values are
ambiguous, but $a-c=s_{12}-s_{23}$ is exactly preserved. The target $a+c$
changes along the missing direction and is not recoverable. Adding
$s_{13}=a+c$ resolves the ambiguity:
\begin{align}
  a&=(s_{12}+s_{13}-s_{23})/2,\nonumber\\
  b&=(s_{12}+s_{23}-s_{13})/2,\nonumber\\
  c&=(s_{13}+s_{23}-s_{12})/2.
  \label{atl:example-inverse}
\end{align}
Known nonzero weights can scale these records without changing their joint
kernel. Collapsing them into one weighted scalar is a further operation that
requires a new observability check. The example supports combining partial
views; it does not assert that arbitrary overlap summaries are sufficient.

For comparison, suppose two slices originally contain $(a,b)$ and $(b,c)$,
but their saved overlap records both contain only $b$. Their agreement operator
is $\delta(e_1,e_2)=e_1-e_2$. Here $\im M=\ker\delta$ holds exactly: every
compatible pair is generated by some source. Nevertheless, $a$ and $c$ remain
unobserved, and $a-c$ cannot be recovered. This example satisfies the universal
compatibility claim while failing target observability. Saving a common
intersection is therefore different from saving enough complementary local
information to reconstruct the desired global distinction.

\subsection{Weighted gluing and target stability}
\label{atl:weighted-gluing}

Given observed records $y$, consider
\begin{equation}
  \widehat z\in\arg\min_z
  \frac12\norm{W^{1/2}(Mz-y)}_2^2,
  \qquad W\succ0.
  \label{atl:weighted-ls}
\end{equation}
The positive matrix $W$ defines the reconstruction metric. Its block weights
can give upper-layer records greater influence. This operation still retains
the indexed record vector $y$ and the map $M$; the scalar objective value is
not their encoding.

\begin{proposition}[Weights and reconstruction error]
\label{atl:weighted-proposition}
Let $H=M^\top WM$, $A=W^{1/2}M$, and
\begin{equation}
  M_W^\dagger=A^\dagger W^{1/2}
  =(M^\top WM)^\dagger M^\top W.
  \label{atl:weighted-inverse}
\end{equation}
Then $\ker H=\ker M$. For $y=Mz+\varepsilon$, the estimate
$\widehat t=CM_W^\dagger y$ satisfies
\begin{equation}
  \widehat t-Cz=-CP_Kz+CM_W^\dagger\varepsilon.
  \label{atl:bias-noise}
\end{equation}
If $\ker M\subseteq\ker C$, the first term vanishes and
\begin{equation}
  \norm{\widehat t-Cz}_2
  \le \norm{CA^\dagger}_2\,
       \norm{W^{1/2}\varepsilon}_2.
  \label{atl:noise-bound}
\end{equation}
\end{proposition}

\begin{proof}
For every $v$, $v^\top Hv=\norm{W^{1/2}Mv}_2^2$. Since $W$ is positive
definite, this expression vanishes exactly when $Mv=0$, proving the kernel
identity. The minimum Euclidean norm least-squares solution is
$A^\dagger W^{1/2}y$. The identity
$A^\dagger=(A^\top A)^\dagger A^\top$ gives
\eqref{atl:weighted-inverse}. Because multiplication by $W^{1/2}$ does not
change the row space, $A^\dagger A=M^\dagger M=I-P_K$. Substituting the
observation model gives
\begin{equation*}
  \widehat t=C(I-P_K)z+CA^\dagger W^{1/2}\varepsilon,
\end{equation*}
which proves \eqref{atl:bias-noise}. Target observability implies $CP_K=0$.
Applying the operator-norm inequality to the remaining term proves
\eqref{atl:noise-bound}.
\end{proof}

Thus positive weights change conditioning and the treatment of inconsistent
observations, but cannot restore a missing target direction. Regularization
can select a solution from an ambiguous family; it does not identify which
member supplied the data. A genuine anchor adds a measurement.

For $r=\rank A>0$ and a compact singular-value decomposition
$A=U_r\Sigma_rV_r^\top$, the target sensitivity factor is
\begin{equation}
  \norm{CA^\dagger}_2
  =\norm{CV_r\Sigma_r^{-1}}_2
  \le \frac{\norm C_2}{\sigma_{\min}^{+}(A)}.
  \label{atl:target-conditioning}
\end{equation}
Small singular values matter through their target coefficients $CV_r$.
If $r=0$, observability requires $C=0$, so the recoverable linear target is
trivial and the singular-value bound is unnecessary.
Moreover, a small gluing residual alone does not bound target error: an error
$\varepsilon=Mv$ is consistent with a different source and can change the
target by $Cv$. The noise bound controls the complete record error, not only
its inconsistent component. All displayed identities use exact arithmetic;
an implementation must also account for rounding and learned-model error.

\subsection{Anchors, gauges, and higher-order distinctions}
\label{atl:gauges}

For a connected graph, let $B$ be its oriented incidence matrix and retain
differences $d=Bz$. If $Bv=0$, every adjacent pair has equal values; connectivity
then makes $v$ constant. Thus $\ker B=\operatorname{span}\{\mathbf1\}$.
Differences determine an absolute state only up to an additive offset. They
determine a target $Cz$ exactly when $C\mathbf1=0$. A positive weighted total
usually fails this test. Retaining one actual node value supplies an anchor;
declaring the mean to be zero merely chooses a convention unless the original
mean was known. For anchor measurements $Gz$, the general condition is
\begin{equation}
  \ker M\cap\ker G\subseteq\ker C.
  \label{atl:anchor-condition}
\end{equation}
Disconnected graphs leave one offset per component; vector-valued states leave
a vector offset in each component.

Higher-order targets can expose ambiguities invisible to pairwise records.
For $x_i\in\{-1,+1\}$, retain all products
$r_{12}=x_1x_2$, $r_{23}=x_2x_3$, and $r_{13}=x_1x_3$.
Both $(1,1,1)$ and $(-1,-1,-1)$ yield the same records and satisfy their
compatibility relation $r_{12}r_{23}r_{13}=1$. Yet the triple product
$x_1x_2x_3$ changes sign. A sign anchor resolves this particular ambiguity.
Compatibility must therefore be tested against the complete target family,
including interactions that become relevant later.

\subsection{What a Hodge decomposition contributes}
\label{atl:hodge}

For an unweighted finite complex, let $B$ take vertex values to edge
differences and $K$ take edge values to oriented face sums. Then $KB=0$.
The edge space has the orthogonal decomposition
\begin{equation}
  E=\im B\ \oplus\ \im K^\top\ \oplus\ \mathcal H,
  \qquad
  \mathcal H=\ker K\cap\ker B^\top.
  \label{atl:hodge-decomposition}
\end{equation}
Indeed, $KB=0$ makes the first two images orthogonal, and their joint
orthogonal complement is precisely the displayed intersection. Also,
\begin{equation}
  \mathcal H=\ker(BB^\top+K^\top K),
  \label{atl:harmonic-kernel}
\end{equation}
because its quadratic form is
$\norm{B^\top e}_2^2+\norm{Ke}_2^2$. Consequently,
$\ker K=\im B\oplus\mathcal H$: locally curl-free records need not be
globally integrable when a harmonic obstruction remains. This distinction,
and reconstruction from pairwise differences, is standard in combinatorial
Hodge theory \citep{hodge}.

Keeping all three components preserves the edge record. Keeping only the
gradient component is a projection. It is justified for exact source recovery
when the allowed records are known to lie in $\im B$, or for a specified task
when that task ignores every discarded ambiguity. Hodge structure can expose
and organize disagreement; it does not establish that curl or harmonic
components are irrelevant to a neural target.

\begin{keybox}
\textbf{The reconstruction contract.}
Save labeled local records before the required distinctions disappear.
Verify their overlap relations, and require that their remaining source
ambiguities are target-invisible. In the linear model these are the separate
conditions $\delta M=0$ and $\ker M\subseteq\ker C$; a universal gluing claim
additionally needs $\im M=\ker\delta$. Use the target sensitivity bound to
assess numerical reliability. These guarantees must be checked again whenever
records are merged, compressed, clipped, or discarded.
\end{keybox}

For nonlinear learned maps, the corresponding global requirement is that equal
complete records imply equal targets. A Jacobian kernel test is only local:
$x\mapsto x^2$ has nonzero derivative on separate positive and negative
intervals while merging their opposite points. Branch identities and
cross-region collisions therefore matter. The linear constructions above
provide controlled cases and explicit certificates; extending them to the
learned nonlinear architecture remains a stated research obligation.
